\chapter{Pleural Fluid Analysis}
\section*{What Is This Test?}
Pleural fluid analysis evaluates fluid aspirated from the pleural space (thoracentesis) to determine the cause of a pleural effusion. The first and most important step is to classify the effusion as a transudate (due to systemic factors affecting the balance of hydrostatic and oncotic pressures across the pleural capillaries) or an exudate (due to local pleural pathology with increased capillary permeability or impaired lymphatic drainage). Light's criteria are used to distinguish transudates from exudates. A pleural effusion is an exudate if it meets ANY ONE of the following: (1) pleural fluid protein/serum protein ratio >0.5, (2) pleural fluid LDH/serum LDH ratio >0.6, or (3) pleural fluid LDH >2/3 of the upper limit of normal for serum LDH. All others are transudates \cite{light1972pleural}. Transudates: congestive heart failure (most common cause overall), cirrhosis (hepatic hydrothorax), nephrotic syndrome, atelectasis, and hypoalbuminemia. Exudates: parapneumonic effusion/empyema (most common exudate), malignancy (lung cancer, breast cancer, lymphoma --- most common causes of malignant effusion), pulmonary embolism, tuberculosis (TB) pleurisy, rheumatoid arthritis, systemic lupus erythematosus, pancreatitis, post-cardiac injury syndrome (Dressler syndrome), drug-induced, and chylothorax.
\section*{Alternative Names}
Thoracentesis Fluid; Pleural Fluid Chemistry; Pleural Effusion Analysis; Light's Criteria
\section*{Principle}
Pleural fluid is collected by thoracentesis under ultrasound guidance. The analysis includes: cell count with differential, total protein, LDH, glucose, pH (collected anaerobically in a heparinized blood gas syringe, placed on ice, analyzed within 30--60 minutes), Gram stain and aerobic/anaerobic culture, cytology for malignant cells, and if indicated: triglycerides (chylothorax, >110 mg/dL), amylase (pancreatic pleural effusion or esophageal rupture), adenosine deaminase/ADA (TB pleurisy, >40--70 U/L), and cholesterol \cite{light2013pleural}.
\section*{History}
Thoracentesis, the aspiration of fluid from the pleural space, has been practiced since antiquity, with descriptions dating to Hippocrates, and was a routine clinical maneuver by the nineteenth century. In 1972, Richard Light and colleagues published the landmark criteria that still classify pleural effusions as transudates or exudates using simultaneously measured pleural fluid and serum protein and LDH. In 1988, Steven Sahn published a widely cited state-of-the-art review of pleural fluid chemistry that systematized the differential diagnosis of the effusion \cite{sahn1988pleura}. The British Thoracic Society issued its first guideline on the investigation of unilateral pleural effusions in adults in 2003, updated in 2010 \cite{maskell2003bts,hooper2010investigation}. Real-time ultrasound guidance, which entered routine practice in the 1990s and early 2000s, markedly reduced the rate of pneumothorax and is now the standard of care for thoracentesis.
\section*{Why This Test Is Ordered}
\begin{itemize}
  \item Evaluation of a newly detected pleural effusion on chest radiograph, ultrasound, or CT to determine whether it is a transudate or exudate and to identify its cause
  \item Classification of a unilateral or unexplained pleural effusion using Light's criteria to direct the search for the underlying disease
  \item Suspected parapneumonic effusion or empyema: pleural fluid pH, glucose, LDH, and cell counts guide the need for chest tube drainage
  \item Suspected malignant effusion: cytologic examination of pleural fluid in patients with known or suspected lung, breast, or other cancer
  \item Suspected tuberculous pleurisy, especially in endemic areas or immunocompromised patients: pleural fluid adenosine deaminase (ADA), AFB smear, and culture
  \item Suspected chylothorax (post-surgical, traumatic, or neoplastic): pleural fluid triglycerides and lipoprotein analysis
  \item Suspected pancreatic pleural effusion or esophageal rupture: pleural fluid amylase
  \item Symptomatic relief: therapeutic thoracentesis of large effusions causing dyspnea
\end{itemize}
\section*{Is This a Primary or Secondary Test}
Pleural fluid analysis is a secondary, confirmatory test that follows detection of a pleural effusion on imaging. The primary diagnostic step is chest radiography or ultrasound demonstrating the effusion, after which thoracentesis and pleural fluid analysis classify the fluid and identify the etiology. In suspected empyema or malignancy, pleural fluid analysis is an essential component of the diagnostic workup.
\section*{How This Test Is Performed}
Thoracentesis is performed with the patient seated upright and leaning forward, or in the lateral decubitus position. Under real-time ultrasound guidance, the skin is cleansed with an antiseptic and a local anesthetic (lidocaine) is infiltrated down to the parietal pleura. A needle or catheter is inserted through the intercostal space just above the rib (to avoid the intercostal neurovascular bundle) into the pleural space, and fluid is aspirated for laboratory analysis. The fluid is distributed into appropriate tubes: a heparinized blood gas syringe for pH (kept on ice and analyzed within 30--60 minutes), EDTA tubes for cell count and differential, plain tubes for chemistry (total protein, LDH, glucose), and culture bottles for aerobic and anaerobic bacteria. The specimen is transported to the laboratory immediately for analysis.
\section*{What Preparation Is Needed}
No fasting or special preparation is required for diagnostic thoracentesis. The patient should sit comfortably and remain still during the procedure. Anticoagulant and antiplatelet medications (warfarin, heparin, direct oral anticoagulants, clopidogrel) should be reviewed with the treating physician, because many current guidelines allow them to be continued for diagnostic thoracentesis but prefer them held for therapeutic taps or when a bleeding diathesis is present. Coagulation studies (INR, platelet count) are often checked beforehand when abnormal bleeding risk is suspected.
\section*{Contraindications and Precautions}
There are no absolute contraindications to thoracentesis when the procedure is clinically necessary. Relative contraindications include an uncorrected bleeding diathesis (coagulopathy, thrombocytopenia, therapeutic anticoagulation), a small or loculated effusion, cellulitis over the puncture site, severe cough or inability to cooperate, and hemodynamic instability. Ultrasound guidance reduces the risk of a dry tap and of pneumothorax. The operator must never advance a needle above the rib margin and should remove only a limited volume (usually 1--1.5 L) during therapeutic thoracentesis to avoid re-expansion pulmonary edema.
\section*{Risks}
Pneumothorax occurs in approximately 1--6\% of procedures with ultrasound guidance and is higher without it. Bleeding (hemothorax, chest wall hematoma), pain, cough, vasovagal reaction, and infection are uncommon, and inadvertent puncture of the liver, spleen, or diaphragm can complicate taps of small or loculated effusions. Re-expansion pulmonary edema is a rare but serious complication of large-volume therapeutic drainage and is minimized by limiting removal to 1--1.5 L per session.
\section*{What Happens After the Test}
The puncture site is covered with a sterile dressing, and a chest radiograph is often obtained after the procedure to exclude pneumothorax. Laboratory results, including the Light's criteria classification and cell counts, are typically available within hours. If an exudate, empyema, or malignancy is identified, the findings direct further evaluation with pleural fluid cytology, closed pleural biopsy, or thoracoscopy, and consultation with pulmonology or thoracic surgery.
\section*{Understanding Results}
\subsection*{Transudate (Normal Pleura, Systemic Cause)}
Congestive heart failure (most common transudate, bilateral effusions, right > left), cirrhosis, nephrotic syndrome, hypoalbuminemia, atelectasis, and peritoneal dialysis.
\subsection*{Exudate (Diseased Pleura)}
Parapneumonic effusion/empyema: WBC markedly elevated (predominantly PMNs), glucose very low (<60 mg/dL), LDH >1,000 IU/L, pH <7.20 (indication for chest tube drainage). Malignancy: bloody or serosanguinous, lymphocyte predominant, cytology positive in 60--70\%, pleural fluid tumor markers (CEA, CA 15-3) may be elevated. TB pleurisy: lymphocyte predominant (>80--90\%), ADA >40--70 U/L (high sensitivity and specificity in high-prevalence areas), AFB smear rare (10--20\%), culture positive in 30--50\%. Rheumatoid arthritis: very low glucose (<30 mg/dL), low pH, elevated LDH, cholesterol crystals (pseudochylous). Chylothorax: milky appearance, triglycerides >110 mg/dL, presence of chylomicrons on lipoprotein electrophoresis.
\section*{Normal Results}
Transudate (any ONE of Light's criteria NOT met). WBC <1,000/mcL. Mononuclear predominance. No malignant cells. Gram stain and culture negative. Glucose approximately equal to serum.
\section*{Follow-up Tests}
Pleural fluid cytology, closed pleural biopsy (Abrams or Cope needle), thoracoscopy with directed pleural biopsy (most sensitive for TB and malignancy), chest CT with contrast, and serum LDH/protein (simultaneously with pleural fluid for Light's criteria).
\section*{References}
\bibliographystyle{unsrtnat}
\bibliography{references}

\clearpage
